\documentclass[11pt,a4paper]{article}
\usepackage[T1]{fontenc}
\usepackage{lmodern}
\usepackage[margin=27mm]{geometry}
\usepackage{amsmath,amssymb,amsthm,mathtools}
\usepackage{microtype}
\usepackage[colorlinks=true,linkcolor=blue,citecolor=blue,urlcolor=blue]{hyperref}
\newtheorem{theorem}{Theorem}[section]
\theoremstyle{remark}
\newtheorem{remark}[theorem]{Remark}
\setlength{\parskip}{4pt}
\setlength{\parindent}{0pt}
\setlength{\emergencystretch}{2em}
\title{Positive kernels and projection couplings:\\two transfers to pure mathematics}
\author{}
\date{Research working note --- 15 November 2026}
\begin{document}
\maketitle
\begin{abstract}
We prove that for every fixed Stirling truncation order $n\ge2$, the
excess of the completely monotonic degree of the signed $m$th derivative
above $2n+m-2$ tends to zero as $m\to\infty$. A quantitative estimate
is $\delta_{n,m}<16/\log(2n+m-1)$ whenever $2n+m-1\ge e^{32}$,
uniformly in $n\ge2$. The proof constructs a negative fractional
Laplace density by comparing a scaled rational kernel with a cosine.
Separately, a monotone coupling of projection determinantal laws
proves the Marcus--de Oliveira enclosure when one matrix has three
eigenvalues, one simple, with the other normal matrix unrestricted.
Both arguments are analytic. Exact finite Stirling degrees and the
general determinantal conjecture remain unresolved here.
\end{abstract}

\paragraph{Status and origin.}
This note records proofs developed by transferring two structures from
work on materials problems: positivity of integral representations,
and orthogonal projections with positive probabilistic realizations.
The first result addresses the fixed-$n$, high-derivative direction
left open in Section~6 of \cite{krasniqi}. The second is a proved
subclass of a longstanding conjecture. The proofs have been checked
within this research session; neither external peer review nor
historical priority is claimed. In particular, a literature search
that finds no exact earlier statement is not a proof of novelty.

\section{The fixed-truncation limit for Stirling remainder degrees}

For integers $n\geq 1$ and real $x>0$, define the signed Stirling remainder by
\[
 R_n(x)=(-1)^n\left[\log\Gamma(x)-\left(x-\frac12\right)\log x+x
 -\frac12\log(2\pi)-\sum_{j=1}^{n}
 \frac{B_{2j}}{2j(2j-1)x^{2j-1}}\right].
\]
Here $B_j$ are the Bernoulli numbers. A smooth function $f$ on $(0,\infty)$
is completely monotone, abbreviated CM, if $(-1)^j f^{(j)}\geq0$ for every
integer $j\geq0$. For the functions below, set
\[
 d_{\rm CM}(f)=\sup\{\rho\geq0:x^\rho f(x)\text{ is CM}\}.
\]
For integers $n\geq2$ and $m\geq0$, write
\[
 f_{n,m}=(-1)^mR_n^{(m)},\qquad p=2n+m-1,\qquad d=p-1,
 \qquad \delta_{n,m}=d_{\rm CM}(f_{n,m})-d.
\]
Krasniqi's preprint, \emph{Noninteger completely monotonic degrees of
Stirling remainders and their derivatives}, arXiv:2610.33699v1 (2026),
proves $0<\delta_{n,m}<1$. Its Section~6 explicitly leaves the behavior
as $m\to\infty$ with $n$ fixed unsettled. The following result addresses
that limiting question. The logarithm in the bound is natural.

\begin{theorem}\label{thm:heat}
For every integer $n\geq2$,
\[
 \boxed{\lim_{m\to\infty}\delta_{n,m}=0.}
\]
More explicitly, uniformly for integers $n\geq2$ and $m\geq0$ satisfying
$p=2n+m-1\geq e^{32}$,
\[
 \boxed{0\leq\delta_{n,m}<\frac{16}{\log p}.}
\]
The constant and threshold are deliberately loose. In particular, this
does not determine the exact degree, a sharp rate, or an asymptotic
constant.
\end{theorem}

\subsection{A Laplace-density criterion}

We include the representation underlying the proof; it is also
Krasniqi's Lemma~3.1 and Proposition~3.2. Put $a_k=2\pi k$ and
\[
 h_p(s)=1-\operatorname{Re}(1+is)^{-p},\qquad s\geq0.
\]
Binet's formula and the partial-fraction expansion of the hyperbolic
cotangent give
\[
 R_n(x)=\int_0^\infty e^{-xt}r_n(t)\,dt,
 \qquad
 r_n(t)=2\sum_{k=1}^\infty
 \frac{t^{2n}}{a_k^{2n}(t^2+a_k^2)}.
\tag{H1}
\]
For completeness, the starting identity is
\[
 \frac1t\left(\frac1{e^t-1}-\frac1t+\frac12\right)
 =2\sum_{k=1}^\infty\frac1{t^2+a_k^2}.
\]
Expanding each summand through $n$ terms by the finite geometric identity
identifies the subtracted coefficients in the definition of $R_n$
and gives (H1).

The elementary identity
\[
 \left(\frac{s^{p+1}}{1+s^2}\right)^{(p-1)}
 =(p-1)!h_p(s)
\]
follows by induction from
$h_{p+1}(s)=h_p(s)+s h'_p(s)/p$.
Consequently,
\[
 \big(t^m r_n(t)\big)^{(d)}
 =H_{n,m}(t):=2d!\sum_{k=1}^\infty a_k^{-2n}h_p(t/a_k),
 \qquad
 x^df_{n,m}(x)=\int_0^\infty e^{-xt}H_{n,m}(t)\,dt.
\tag{H2}
\]
The series and the derivatives used here converge uniformly on compact
intervals. At zero, $t^m r_n(t)=O(t^{d+2})$, so the boundary terms at
zero in the $d$ integrations by parts vanish. Since $|h_p|\leq2$, the
right-hand side defining $H_{n,m}$ is bounded, and integration gives
polynomial growth of the lower derivatives; the boundary terms at
infinity therefore vanish as well. Also $h_p(s)\geq0$, which shows
that $x^df_{n,m}$ is CM and hence $\delta_{n,m}\geq0$.

For $0<\alpha<1$ and a differentiable function $h$ with $h(0)=0$, define
\[
 D^\alpha h(t)=\frac1{\Gamma(1-\alpha)}
 \int_0^t(t-u)^{-\alpha}h'(u)\,du.
\tag{H3}
\]
This is both the Caputo and Riemann--Liouville derivative under this
zero initial condition. Its Laplace transform is $x^\alpha\mathcal Lh(x)$.
Thus, with
\[
 Q_{n,m,\alpha}(t)
 =2d!\sum_{k=1}^\infty
 a_k^{-2n-\alpha}(D^\alpha h_p)(t/a_k),
\]
we have
\[
 x^{d+\alpha}f_{n,m}(x)=\mathcal L Q_{n,m,\alpha}(x),
 \qquad
 x^{d+\alpha}f_{n,m}\text{ is CM}
 \quad\Longleftrightarrow\quad
 Q_{n,m,\alpha}(t)\geq0\ \text{for all }t>0.
\tag{H4}
\]
Indeed, $|h'_p|\leq p$ yields
$|D^\alpha h_p(t)|\leq pt^{1-\alpha}/\Gamma(2-\alpha)$.
The series for $Q$ therefore converges locally uniformly, is continuous,
and has at most polynomial growth. These bounds also justify the Laplace
interchanges. For the converse implication in (H4), Bernstein's theorem
represents a CM function as the Laplace transform of a nonnegative
measure. Uniqueness of the Laplace transform, after exponentially
weighting to obtain finite measures, identifies that measure with
$Q(t)\,dt$. Continuity then gives pointwise nonnegativity.

Let
\[
 g_p(u)=h_p(u/p),\qquad g(u)=1-\cos u,
 \qquad
 S_{p,\alpha}(y)=\sum_{k=1}^\infty
 k^{-2n-\alpha}(D^\alpha g_p)(y/k),
\]
and define $S_{\infty,\alpha}$ by replacing $g_p$ with $g$.
The scaling in (H3) gives the exact normalization
\[
 Q_{n,m,\alpha}(2\pi y/p)
 =2d!\,(2\pi)^{-2n-\alpha}p^\alpha S_{p,\alpha}(y).
\tag{H5}
\]
It will therefore suffice to find a negative value of $S_{p,\alpha}$.

\subsection{An oscillatory fractional derivative}

Write $D_\alpha(u)=D^\alpha(1-\cos u)$. Substituting $v=u-t$ in
(H3) gives
\[
 D_\alpha(u)=\frac1{\Gamma(1-\alpha)}
 \operatorname{Im}\left[e^{iu}\int_0^u v^{-\alpha}e^{-iv}\,dv\right].
\]
The improper integral over $(0,\infty)$ is
$\Gamma(1-\alpha)e^{-i\pi(1-\alpha)/2}$, by the usual Abel limit of
the gamma integral. Integration by parts on the tail gives
\[
 \left|\int_u^\infty v^{-\alpha}e^{-iv}\,dv\right|
 \leq 2u^{-\alpha}.
\]
It follows that, for every $u>0$,
\[
 \left|D_\alpha(u)+\cos(u+\pi\alpha/2)\right|
 \leq\frac{2u^{-\alpha}}{\Gamma(1-\alpha)}.
\tag{H6}
\]
Splitting the same integral at $1$ and integrating the part over
$[1,u]$ by parts also gives the global bound
\[
 |D_\alpha(u)|\leq
 \frac{(1-\alpha)^{-1}+2}{\Gamma(1-\alpha)}.
\tag{H7}
\]
For $0<s\leq1$, Jensen's inequality for an exponential random variable
of mean $1$ gives $\Gamma(s)\geq1$. Hence, for $0<\alpha\leq1/2$,
\[
 |D_\alpha(u)|\leq4,
 \qquad \Gamma(2-\alpha)\geq\tfrac12.
\tag{H8}
\]

\subsection{A negative density with explicit parameters}

Assume $p\geq e^{32}$ and choose
\[
 \alpha=\frac{16}{\log p},\qquad
 y=2\pi\left\lfloor\frac{p^{1/4}}{2\pi}\right\rfloor.
\tag{H9}
\]
Then $0<\alpha\leq1/2$, $y\in2\pi\mathbb N$, and
\[
 \frac12p^{1/4}\leq y\leq p^{1/4},\qquad
 y^{-\alpha}\leq\sqrt2e^{-4}.
\]
Using (H6) for the first summand and (H8) for the rest,
\begin{align*}
 S_{\infty,\alpha}(y)
 &\leq-\cos(\pi\alpha/2)
       +\frac{2y^{-\alpha}}{\Gamma(1-\alpha)}
       +4\sum_{k=2}^\infty k^{-2n-\alpha}\\
 &\leq-\frac1{\sqrt2}+2\sqrt2e^{-4}+\frac5{12}.
\end{align*}
Here we used $n\geq2$ and
\[
 \sum_{k=2}^\infty k^{-4}
 \leq\frac1{16}+\int_2^\infty t^{-4}\,dt=\frac5{48}.
\tag{H10}
\]

We next estimate the difference between the finite-$p$ and limiting
kernels on the growing interval $[0,y]$. Explicit differentiation gives
\[
 g'_p(u)=A_p(u)\sin\theta_p(u),\qquad
 A_p(u)=\left(1+\frac{u^2}{p^2}\right)^{-(p+1)/2},
 \qquad \theta_p(u)=(p+1)\arctan(u/p).
\]
For $0\leq u\leq y\leq p$ and $p\geq1$,
\[
 1-A_p(u)\leq\frac{u^2}{p},\qquad
 |\theta_p(u)-u|
 \leq\frac up+\frac{2u^3}{3p^2}
 \leq\frac up+\frac{2u^2}{3p}.
\]
The first inequality follows from $1-e^{-z}\leq z$ and
$\log(1+z)\leq z$. The second follows from
$0\leq v-\arctan v\leq v^3/3$.
As $y\geq1$, this implies
\[
 \sup_{0\leq u\leq y}|g'_p(u)-\sin u|
 \leq\frac{8y^2}{3p}.
\tag{H11}
\]
Applying (H3) at $y/k$ and summing therefore yields
\begin{align*}
 |S_{p,\alpha}(y)-S_{\infty,\alpha}(y)|
 &\leq\frac{8y^{3-\alpha}}{3p\,\Gamma(2-\alpha)}
         \sum_{k=1}^\infty k^{-2n-1}\\
 &\leq\frac{67}{12}p^{-1/4}.
\tag{H12}
\end{align*}
In the last step, $y\geq1$, $y\leq p^{1/4}$, (H8), and the bound
\[
 \sum_{k=1}^\infty k^{-5}
 \leq1+\frac1{32}+\int_2^\infty t^{-5}\,dt=\frac{67}{64}
\]
were used. Combining the estimates gives
\[
 S_{p,\alpha}(y)
 \leq-\frac1{\sqrt2}+2\sqrt2e^{-4}
       +\frac5{12}+\frac{67}{12}e^{-8}<0.
\tag{H13}
\]
The sign can be checked entirely with rational bounds:
$1/\sqrt2>7/10$, $2\sqrt2<3$, $e^4>50$, and $e^8>2500$
give an upper bound of $-2211/10000$. The exponential bounds follow,
for example, from $e>8/3$.

By (H5), the continuous density $Q_{n,m,\alpha}$ is negative at
$t=2\pi y/p$. Consequently $x^{d+\alpha}f_{n,m}$ is not CM by (H4).
The set of admissible exponents is downward closed: multiplying a CM
function by $x^{-b}$, $b\geq0$, preserves complete monotonicity.
It is also closed, since each derivative of $x^\rho f_{n,m}(x)$ depends
continuously on $\rho$ at each fixed $x>0$.
The failure at $d+\alpha$ thus gives
$d_{\rm CM}(f_{n,m})<d+\alpha$.
Together with $\delta_{n,m}\geq0$, this proves the explicit bound and,
for every fixed $n\geq2$, the asserted limit.\hfill$\square$

\paragraph{An alternative qualitative argument.}
The convergence alone has a shorter proof, independent of the explicit
parameter choices. For fixed $0<\alpha<1$, one has
$g_p\to1-\cos$ in $C^1$ on compact intervals and $|g'_p|\leq1$.
Thus dominated convergence gives
$S_{p,\alpha}(y)\to S_{\infty,\alpha}(y)$ for every fixed $y$, since
\[
 |D^\alpha g_p(y/k)|\leq
 \frac{(y/k)^{1-\alpha}}{\Gamma(2-\alpha)}.
\]
Let $y_N=2\pi N!$. For each fixed $k$, $y_N/k$ is an integer multiple
of $2\pi$ for all sufficiently large $N$, and tends to infinity.
Equations (H6)--(H7) and dominated convergence imply
\[
 S_{\infty,\alpha}(y_N)
 \longrightarrow
 -\cos(\pi\alpha/2)\sum_{k=1}^\infty k^{-2n-\alpha}<0.
\]
Choose $N$ first so that the limiting density is negative, and then
let $p$ grow at that fixed $y_N$. This gives $\delta_{n,m}<\alpha$
eventually. Since $\alpha>0$ is arbitrary, the limit follows.

\section{A determinantal enclosure from a nested coupling}

The Marcus--de Oliveira conjecture asserts that for normal matrices
$A,B\in\mathbb C^{N\times N}$ with eigenvalues $a_i,b_i$,
\begin{equation}\label{eq:moc}
 \det(A+B)\in\operatorname{conv}
 \left\{\prod_{i=1}^N(a_i+b_{\sigma(i)}):\sigma\in S_N\right\}.
\end{equation}
The convex hull is taken in $\mathbb C\cong\mathbb R^2$.
The general conjecture remains outside the scope of this note;
\cite{survey} reviews its known special cases.

\begin{theorem}\label{thm:moc}
Let $A$ be any normal $N\times N$ matrix. Suppose
\[
 B=\beta I+bP+cQ,\qquad Q=vv^*,\quad \|v\|_2=1,\quad Pv=0,
\]
where $P$ is an orthogonal projection and $\beta,b,c\in\mathbb C$.
Then \eqref{eq:moc} holds. In particular, it holds whenever $B$
has three distinct eigenvalues, one of them simple, with arbitrary
multiplicities for the other two.
\end{theorem}

\begin{proof}
After a simultaneous unitary change of coordinates, write
$D=\operatorname{diag}(d_1,\ldots,d_N)$ with $d_i=a_i+\beta$.
We must enclose $\det(D+bP+cQ)$.

For an orthogonal projection $R$, its determinantal probability law
is a random subset $U\subseteq[N]$ satisfying
\begin{equation}\label{eq:inclusion}
 \mathbb P(I\subseteq U)=\det R[I,I],\qquad |U|=\operatorname{rank}R.
\end{equation}
These facts hold over the complex Hilbert space as well as the real one.
Lyons's stochastic domination theorem \cite[Theorem 6.2]{lyons},
together with the monotone-coupling characterization of stochastic
domination, provides a joint law
\[
 S\sim\mathbb P_P,\qquad T\sim\mathbb P_{P+Q},\qquad S\subseteq T.
\]
Here $P+Q$ is an orthogonal projection. If $r=\operatorname{rank}P$,
then $|S|=r$, $|T|=r+1$, so $T\setminus S=\{j\}$ almost surely.

For a subset $U$ put
\[
 F_U(b)=\prod_{i\in U}(d_i+b)\prod_{i\notin U}d_i.
\]
Expanding the product and using \eqref{eq:inclusion} gives
\begin{align}
 \mathbb E F_S(b)
 &=\sum_{I\subseteq[N]} b^{|I|}\det P[I,I]\prod_{i\notin I}d_i
 =\det(D+bP),\label{eq:principalexp}\\
 \mathbb E F_T(b)&=\det(D+b(P+Q)).\nonumber
\end{align}
No nonzero-diagonal or invertibility assumption is used.

Define
\[
 Z=\prod_{i\in S}(d_i+b)(d_j+c)
       \prod_{i\notin S\cup\{j\}}d_i.
\]
Every value of $Z$ is a vertex appearing on the right side of
\eqref{eq:moc}: it assigns $r$ copies of $\beta+b$, one copy of
$\beta+c$, and $N-r-1$ copies of $\beta$ to the $a_i$.
For $b\ne0$ the following identity holds for every outcome:
\[
 Z=F_S(b)+\frac cb\bigl(F_T(b)-F_S(b)\bigr).
\]
Since $Q$ has rank one, $\det(D+bP+cQ)$ is an affine polynomial in
$c$. Therefore \eqref{eq:principalexp} implies
\begin{align*}
 \mathbb E Z
 &=\det(D+bP)+\frac cb
    \bigl[\det(D+b(P+Q))-\det(D+bP)\bigr]\\
 &=\det(D+bP+cQ).
\end{align*}
Thus the determinant is a convex combination with genuine
nonnegative probability weights. The case $b=0$ follows by
continuity with the same coupling, or directly by the rank-one
principal-minor expansion.
\end{proof}

\begin{remark}[Scope and what the argument does not establish]
The first multiplicity pattern covered here beyond the familiar
``all but two eigenvalues equal'' family is $(2,2,1)$ in dimension
$5$. In every dimension the first matrix $A$ is unrestricted apart
from normality. No claim of historical priority for this subclass
is made: the theorem and its proof are established here, while
exhaustive comparison with prior special-case results remains to
be done. The full conjecture is not proved. If $Q$ has rank greater
than one, the determinant is no longer affine in $c$; the two
coupled marginal laws do not supply the additional mixed moments.
\end{remark}

\begin{remark}[Finite certificate]
The convex weights can be obtained by a finite transportation
problem. Its left nodes are the $r$-subsets $S$ with supplies
$\det P[S,S]$, and its right nodes are the $(r+1)$-subsets $T$
with demands $\det(P+Q)[T,T]$. Connect $S$ to $T$ exactly when
$S\subset T$. The coupling theorem guarantees a nonnegative flow
of total mass one; each edge flow is a coefficient of the
corresponding permutation vertex above.
\end{remark}

\begin{thebibliography}{9}
\bibitem{krasniqi}
V. Krasniqi, \emph{Noninteger completely monotonic degrees of Stirling
remainders and their derivatives}, arXiv:2610.33699v1 (2026).
\href{https://arxiv.org/abs/2610.33699}{arxiv.org/abs/2610.33699}.

\bibitem{lyons}
R. Lyons, \emph{Determinantal probability measures},
Publications Math\'ematiques de l'IH\'ES \textbf{98} (2003), 167--212.
Theorem 6.2 and the monotone-coupling characterization of stochastic
domination.
\href{https://doi.org/10.1007/s10240-003-0016-0}{doi:10.1007/s10240-003-0016-0}.

\bibitem{survey}
N. Bebiano and J. P. da Provid\^encia,
\emph{Revisiting the Marcus--de Oliveira conjecture},
Mathematics \textbf{13} (2025), article 711.
\href{https://doi.org/10.3390/math13050711}{doi:10.3390/math13050711}.
\end{thebibliography}
\end{document}
